\documentclass[10pt]{article}

\usepackage[utf8]{inputenc}

\usepackage[T1]{fontenc}

\usepackage{amsmath}

\usepackage{amsfonts}

\usepackage{amssymb}

\usepackage[version=4]{mhchem}

\usepackage{stmaryrd}

\usepackage{graphicx}

\usepackage[export]{adjustbox}

\graphicspath{ {./images/} }

\usepackage{bbold}



\begin{document}

Для того чтобы функция $f(y), y \in \Gamma$, была разрешимой, необходимо и достаточно, чтобы она была интегрируемой по гармонич. мере на Г (т е о р е м а Б р е ло). Јюбая непрерывная конечная функция $f(y), y \in \Gamma$, разрешима (т е о р е м а



Точка $y_{0} \in \Gamma$ наз. $\mathrm{p}$ е $\mathrm{r} \mathbf{л}$ л я $\mathrm{p}$ н о й $\mathrm{r} \mathrm{p}$ а и ч ч о й то ч к о й, если для любой непрерывной конечной фуунции $f(y), \quad y \in \Gamma$, выполняется предельное соотношение



\[

\lim _{x \rightarrow y_{0}} H_{f}(x)=f\left(y_{0}\right)

\]



Регулярность всех точек $y \in \Gamma$ равносильна существованию классич. решения $w_{f}(x)$ задачи Дирихле для любой непрерывной конечной функции $f(y), y \in \Gamma$, причем в этом случае $H_{f}(x)=w_{f}(x)$; область $\Omega$, все граничные точки к-рой регулярны, иногда наз. также p e г у л я р н о й. Для того чтобы точка $y_{0} \in \Gamma$ была регулярной, необходимо и достаточно, чтобы существовал барьер в $y_{0}$.



Точки $y_{0} \in \Gamma$, не являющиеся регулярными, наз. и р регу ля рным и г раничными точкам и. Напр., иррегулярными граничными точками являются изолированные точки и при $n \geqslant 3$ вершины достаточно сильно заостренных входящих в область $\Omega$ острий (п р п м ө р Л е б е г а). Множество всех иррегулярных точек $\Gamma$ есть множество тина $F_{\sigma}$ емкости нуль.



Пусть имеется последовательность областей $\Omega_{k}$, $\bar{\Omega}_{k} \subset \Omega_{k+1}$, такая, что $\Omega=\bigcup_{k=1}^{\infty} \Omega_{k}, \quad$ и непрерывная конечная функция $f(y), y \in \Gamma$, продолжена непрерывно на окрестность $\Gamma$. Тогда



\[

\lim _{k \rightarrow \infty} H_{f}\left(x ; \Omega_{k}\right)=H_{f}(x ; \Omega), x \in \Omega,

\]



равномерно внутри $\Omega$; в случае регулярных областей $\Omega_{k}$ здесь получается конструкция обобщенного решения зацачи Дирихле по Винеру. Рассмотрим теперь для области $\Omega$ без внутренней границы произвольную последовательность областей $G_{k}, \quad \partial G_{k} \rightarrow \Gamma, G_{k} \supset \bar{\Omega}$. В этом случае, вообще говоря,



\[

\lim _{k \rightarrow \infty} H_{f}\left(x ; G_{k}\right) \neq H_{f}(x ; \Omega) .

\]



Задача Дирихле устойчива в области $\Omega$ или в замкнутой области $\bar{\Omega}$, если



\[

\lim _{k \rightarrow \infty} H_{f}\left(x ; G_{k}\right)=H_{f}(x ; \Omega)

\]



соответственно для всех $x \in \Omega$ или для всех $x \in \bar{\Omega}$. Для устойчивости задачи Дирихле в области $\Omega$ необходимо и достаточно, чтобы множества всех иррегулярных точек дополнений $C \Omega$ и $C \bar{\Omega}$ совпадали; для устойчивости в замкнутой области - чтобы дополнение $C \bar{\Omega}$ не имело иррегулярных точек (т е о р е мы д ыш а, см. [4], где построен также пример регулярной области $\Omega$, внутри к-рой задача Дирихле неустойчива).



См. также Верхних и нижних Функиий метод.



Jum.: [1] P e r r o n O., "Math. Z.", 1923, Bd 18, S. 42-54;

[ ] е т р о в с к и й И. Г., "Успехи матем. наук", 1941 [1940], [2] П е т р о в с к и й И. Г., "Успехи матем. наук", 1941 [1940],



\includegraphics[max width=\textwidth, center]{2023_07_08_2f4fe53faf38228a7281g-1(1)}

д ы ї М. В., "Успехи матем. наук", 1941 p. $[1940]$, в. 8, с. 171231 ; [5] Бр е л о М., Основы классической теории потенциала, пер. с франц., М., 1964.

ПЕРРОНА ПРЕОБРАЗОВАНИЕ - ортогональное (унитарное) преобразование



\[

x^{i}=\sum_{j=1}^{n} u_{j}^{i}(t) y^{j}, i=1, \ldots, n,

\]



гладко зависящее от $t$ и преобразующее линейную систему обыкновенных дифференциальных уравнений



\[

\dot{x}^{i}=\sum_{j=1}^{n} a_{j}^{i}(t) x^{j}, i=1, \ldots, n,

\]



в систему треугольного вида



\[

\dot{y}^{i}=\sum_{j=i}^{n} p_{j}^{i}(t) y^{j}, i=1, \ldots, n .

\]



Введено О. Перроном [1]. Справедлива т е о р е м а П е р р о н а : для всякой линейной системы (2) с непрерывными коэффициентами $a_{j}^{i}(t)$ существует П. п.



П. п. строится с помощью процесса ортогонализации Грама - Шмидта (при каждом $t$ ) системы векторов $x_{1}(t), \ldots, x_{n}(t)$, где $x_{1}(t), \ldots, x_{n}(t)$ - какая-либо бундаментальная система решений системы (2), причем различные фундаментальные системы дают, вообще говоря, различные П. П. (см. [1], [2]). Для систем (2) c ограниченными непрерывными коэффициентами все П. П. являются Ляпунова преобразованияли.



Если матричнозначная функция $\left\|a_{j}^{i}(t)\right\|, i, j=1, \ldots$, $n$, является рекуррентной функцией, то найдется рекуррентная матричнозначная функция $\left\|u_{j}^{i}(t)\right\|, i, j=1$, $\ldots, n$, такая, что (1) есть П. П., приводящее систему (2) к треугольному виду (3), причем функция



\[

\left\|p_{j}^{i}(t)\right\|, i, j=1, \cdots, n,

\]



рекуррентна.



[2] Jum.: [1] P e r r o n O., "Math. Z.", 1930, Bd 32, S. 465-73; 1-38; [3] Б ы ло в Б. Ф., В и но о р а д Р.'Э., ${ }^{19}$ р о б ма н Д. М., Н е м ы ц к и й В. В., Теория показателей Јяпуунова и ее приложения к вопросам устойчивости, М., 1966; [4] И 3 о-

б о в Н. А., в сб.: Итоги науки и техники. Математический ана-



\begin{center}

\includegraphics[max width=\textwidth]{2023_07_08_2f4fe53faf38228a7281g-1}

\end{center}



ПЕРРОНА - СТИЛТЬЕСА ИНТЕГРАЛ - обобщение понятия Перрона интеграла от функции одного действительного переменного. Конечная функция $f(x)$ наз. интегрируемой в смысле Перрона - Стилтьеса относительно конечной функции $G(x)$ на $[a, b]$, если на $[a, b]$ существуют мажоранта $M(x)$ и миноранта $m(x)$ функции $f(x)$ относительно $G(x)$ на $[a, b]$ с $\boldsymbol{M}(x)=$ $=m(a)=0$ такие, что в каждой точке $x \in[a, b]$



\[

M(x+\beta)-M(x-\alpha) \geqslant f(x)(G(x+\beta)-G(x-\alpha)),

\][



m(x+\textbackslash beta)-m(x-\textbackslash alpha) \textbackslash leqslant f(x)(G(x+\textbackslash beta)-G(x-\textbackslash alpha))

]



при всех достаточно малых $\alpha \geqslant 0$ и $\beta \geqslant 0$, кроме того, нижняя грань чисел $M(b)$, где $M(x)$ - произвольная мажоранта $f(x)$ относительно $G(x)$, и верхняя грань чисел $m(b)$, где $m(x)$ - произвольная миноранта $f(x)$ относительно $G(x)$, равны между собой. Общее значение этих двух граней наз. П.-С. и. от $f(x)$ относительно $G(x)$ на $[a, b]$ и обозначается



\[

(P-S) \int_{a}^{b} f(x) d G(x)

\]



[1].



Такое обобщение интеграла Перрона ввел О. Уорд



Jum.: [1] W a r d A. J., "Math. Z.", 1936, Bd 41, S. 578604; [2] с а к с С., Теория интеграла, пер. с англ., М., 1949; [3] В и н о г р а до в а И. А., С к в о р ц о в В. А., В Кн.: Мате-

матический анализ. 1970, М.,1971, с. 65-107. T. П. Лyкашенко.



IЕРРОНА - ФРОБЕНИУСА ТЕОРЕМА: пустЬ действительная квадратная $(n \times n)$-матрица $A$, рассматриваемая как оператор в пространстве $\mathbb{R}^{n}$, не имеет инвариантных координатных подпространств (такая матрица наз. н е р а з л о ж и м о й) и неотрицательна (т. е. все ее элементы неотрицательны). И пусть $\lambda_{1}, \ldots, \lambda_{n}$ - ее собственные значения, занумерованные так, что



\[

\left|\lambda_{1}\right|=\ldots=\left|\lambda_{h}\right|>\left|\lambda_{h+1}\right| \geqslant \ldots \geqslant\left|\lambda_{n}\right|, 1 \leqslant h \leqslant n .

\]



Тогда:



\begin{enumerate}

  \item число $r=\left|\lambda_{1}\right|$ - простой положительный корень характеристич. многочлена матрицы $A$;



  \item существует собственный вектор матрицы $A$ с положительными координатами, соответствующий $r$;



  \item числа $\lambda_{1}, \ldots, \lambda_{h}$ совпадают с точностью до нумерации с числами $r, \omega r, \ldots, \omega^{h-1} r$, где $\omega=e^{2 \pi i / h}$;



\end{enumerate}



\end{document}